\section{Types}
\label{sec:types}

We augment the definition of conditionally linear functions with a construct
we call \emph{types}.
A type~$\tvar$ is an element of a \emph{type set} $\type$, and a $\type$-typed
family of conditionally linear functions is a collection
$\{L_\tvar\}_{\tvar\in\type}$ containing a CL function $L_\tvar$ for each
type~$\tvar\in\type$.
The utility of this definition is that it allows us to define another object,
namely conditionally linear distributions parameterized by an undirected graph
$G = (\type, E)$ on the set of types known as a \emph{type graph}.
Given two $\type$-typed families of conditionally linear functions
$\{L_\lvar\}_{\lvar\in\type}, \{R_\rvar\}_{\rvar\in\type}$, the $(\type,
G)$-typed conditionally linear distribution corresponding to them is the
distribution which samples a pair of types $(\lvar, \rvar)$ uniformly at random
from the edges of~$G$ (with each endpoint having equal probability as being
chosen for $\lvar$ or $\rvar$, respectively) and then samples $(x, y)$ from
$\mu_{L_\lvar,\, R_\rvar}$.
The output is the pair $( (\lvar, x), (\rvar, y) )$.

The normal form verifiers we present in the paper frequently use typed CL
distributions to sample their questions, rather than untyped CL distributions.
Types allow us to model the parts of their question distributions which are
unstructured and unsuitable for being sampled from CL distributions.
A common use of types is to allow the verifier to use previously defined games
as subroutines.
Here, the type helps indicate which subroutine the verifier selects, and an edge
in the type graph between two different types allows us to introduce a test that
cross-checks the results of one subroutine with the results of another.

Finally, we show how to convert any typed CL distribution into an equivalent (in
the precise sense defined below) untyped CL distribution with two additional
levels, a technique we call \emph{detyping}.
This entails showing how to ``simulate'' the \emph{graph distribution} of $G =
(\type, E)$, i.e.\ the uniform distribution on its edges, using an untyped CL
distribution.
The simulation we give is based on rejection sampling and is only approximate:
its quality degrades exponentially with the number of types in~$\type$.
As a result, we will ensure throughout the paper that all type sets we consider
are of a small, in fact generally constant, size.

This section is organized as follows.
In \cref{sec:typed-samplers} we define typed variants of CL distributions,
samplers, deciders, and verifiers.
In \cref{sec:graph-dist} we define a CL distribution which samples from the
graph distribution of a given graph $G = (\type, E)$.
In \cref{sec:detype} we define a canonical way to detype typed samplers,
deciders, and verifiers using the graph sampler from \cref{sec:graph-dist}.
We then prove the main result of the section, \cref{lem:detyping-verifiers},
which relates the value of the detyped normal form verifier to the value of
the original typed verifier.

\subsection{Typed samplers, deciders, and verifiers}
\label{sec:typed-samplers}

\begin{definition}[Typed conditionally linear functions]
  Let $\type$ be a finite set and $V$ be $\F^n$ for some integer $n \geq 0$.
  A \emph{$\type$-typed family of $\ell$-level conditionally linear functions}
  (implicitly, \emph{on $V$}) is a collection $\{L_\tvar\}_{\tvar \in \type}$
  such that, for each $\tvar \in \type$, $L_\tvar$ is an $\ell$-level
  conditionally linear function on~$V$.
\end{definition}

\begin{definition}[Graph distribution]
  \label{def:graph-distribution}
  Let $G = (U, E)$ be an undirected graph with vertex set $U$ and edge set~$E$.
  Edges in~$E$ are written as multisets $\{u, v\}$ of two vertices; the case $u
  = v$ represents a self-loop.
  Suppose there are $m$ edges, $k$ of which are self-loops.
  Then the \emph{graph distribution $\mu_G$ of~$G$} is the distribution over
  $U\times U$ such that for every $(u, v)\in U\times U$,
  \begin{equation*}
    \mu_G(u, v) =
    \begin{cases}
      1/(2m-k) & \text{if $\{u, v\} \in E$},\\
      0 & \text{otherwise}.
    \end{cases}
  \end{equation*}
  This is identical to the uniform distribution over pairs $(u, v)\in U\times U$
  such that $\{u, v\} \in E$.
\end{definition}

\begin{definition}[Typed conditionally linear distributions]
  Let $\type$ be a type set and
  $L=\{L_\lvar\}_{\lvar\in\type},R=\{R_\rvar\}_{\rvar\in\type}$ be $\type$-typed
  families of conditionally linear functions on~$V$.
  Let $G = (\type, E)$ be a graph with vertex set $\type$.
  The \emph{$(\type, G)$-typed conditionally linear distribution $\mu_{L,\,
      R}^G$ corresponding to~$(L, R)$} is the distribution over pairs $( (\lvar,
  x), (\rvar, y) )$, where $(\lvar, \rvar)$ is drawn from $\mu_G$ and $(x,y)$ is
  drawn from $\mu_{L_\lvar,\, R_\rvar}$.
\end{definition}

\begin{definition}[Typed conditionally linear samplers]
  \label{def:typed-sampler}
	Let $q:\N \to \N$ be an admissible field size function and $s: \N \to \N$ be a
  function.
  Let $\type$ be a finite type set and let $G = (\type, E)$ be a graph with vertex set $\type$.
  A $7$-input Turing machine $\sampler$ is a \emph{$(\type, G)$-typed, $\ell$-level
    conditionally linear sampler with field size $q(n)$ and dimension $s(n)$} if
  for all $n \in \N$, letting $q = q(n)$ and $s = s(n)$, there exist
  $\type$-typed families of $\ell$-level conditionally linear functions
  $\{L^{\alice,\, n}_{\tvar}\}_{\tvar \in \type}$ and $\{L^{\bob,\,
    n}_{\tvar}\}_{\tvar \in \type}$ on $V = \F_q^s$ where $\tvar \in \type, w\in
  \AB$, the conditionally linear function $L^{w,\, n}_{\tvar}$ has marginal
  functions $\{L^{w,\, n}_{\tvar,\, \leq j}\}$ and factor spaces $\{V^{w,\,
    n}_{\tvar,\, j,\, u}\}$ satisfying the conditions of Lemma~\ref{lem:cl-kth},
  and for all $\tvar \in \type$, $w \in \ab$, $j \in \{1, \ldots, \ell\}$, and
  $z \in V$:
  \begin{itemize}
  \item On input $(n, \gamestyle{dimension})$, the sampler $\sampler$ returns
    the dimension $s(n)$.
	
	\item On input $(n, w, \gamestyle{marginal}, j, z, \tvar)$, the sampler
    $\sampler$ returns the binary representation of $L^{w,\, n}_{\tvar,\, \le
      j}(z)$.
			
			\item On input $(n, w, \gamestyle{linear}, j, u, y, \tvar)$, the sampler $\sampler$
    outputs the binary representation of $L^{w,\, n}_{\tvar,\, j,\, u}(y)$,

	\item On input $(n, w, \gamestyle{factor}, j, u, \tvar)$, the sampler
    $\sampler$ returns the factor space $V^{w,\, n}_{\tvar,\, j,\, u}$ of
    $L^{w,\, n}_{\tvar}$, represented as an indicator vector in $\{0,1\}^s$.
  \end{itemize}
  We call $\F_{q(n)}^{s(n)}$ the \emph{ambient space of $\sampler$}.
  We call $\{ L^{\alice,\, n}_{\tvar} \}, \{ L^{\bob,\, n}_{\tvar} \}$ the
  \emph{CL functions of $\sampler$ on index $n$}.
  The \emph{time complexity} of $\sampler$, denoted $\TIME_\sampler(n)$, is the
  number of steps before $\sampler$ halts for index $n$.
\end{definition}

We assume that types $\tvar \in \type$ are represented using binary strings of
length at most $\lceil \log |\type| \rceil$; if a type $\tvar$ is given as input
to the sampler $\sampler$ and is not an element of $\type$, then the sampler
returns $0$.
Furthermore, as described in Remark~\ref{rmk:sampler-inputs} for un-typed
samplers, we write typed samplers with different numbers of arguments depending
on the input.
We note that Definition~\ref{def:typed-sampler} includes the graph $G$ as a parameter, even though it is not explicitly used anywhere;
this is so we can define the following concept, i.e.\ the \emph{distribution} of a typed sampler, which \emph{does} depend on $G$.

\begin{definition}[Distribution of a typed sampler]
  \label{def:typed-sampler-sample}
  Let $\sampler$ be a $(\type, G)$-typed sampler.
  Let $L^w = \{ L^{w}_{\tvar} \}_\tvar$ for $w\in \{\alice, \bob\}$ be the CL
  functions of $\sampler$ on index~$n$.
  The \emph{distribution of sampler $\sampler$ on index $n$},
  denoted $\mu^G_{\sampler,\, n}$, is the $(\type, G)$-typed conditionally
  linear distribution corresponding to $(L^\alice, L^\bob)$.
\end{definition}

\begin{definition}[Downsizing typed CL samplers]
  \label{def:downsize-typed-sampler}
Let $\sampler$ be a typed sampler.
  The typed downsized sampler $\downsize(\sampler)$ is defined as in
  Definition~\ref{def:downsize_sampler} with the only difference that the type
  $\tvar$ is included as part of the input to the sampler, as in
  Definition~\ref{def:typed-sampler}.
\end{definition}

\begin{lemma}
  \label{lem:downsize_typed_sampler}
Let $\sampler$ be a $(\type, G)$-typed $\ell$-level CL sampler, for some finite set
  $\type$, graph $G$,  and integer $\ell\geq 0$.
  Let $q(n)$ and $s(n)$ be as in Definition~\ref{def:typed-sampler}.
  Then $\downsize(\sampler)$ defined in
  Definition~\ref{def:downsize-typed-sampler} is a $(\type, G)$-typed $\ell$-level CL
  sampler with field size $2$, dimension $s(n) \log q(n)$, and
  time complexity
  \begin{equation*}
    \TIME_{\downsize(\sampler)}(n) = O(\TIME_\sampler(n) \log q(n))\;.
  \end{equation*}
  Furthermore, for every integer $n \geq 1$, the CL functions of
  $\downsize(\sampler)$ on index $n$ are
  $\{(L_\tvar^{w,n})^\downsize\}_{w\in\{\alice,\bob\}, \tvar\in\type}$, as
  defined in Definition~\ref{def:cl-downsize}.
\end{lemma}

\begin{proof}
  The proof is analogous to the proof of Lemma~\ref{lem:downsize_sampler}, and
  we omit it.
\end{proof}

\begin{definition}[Typed decider]
  \label{def:typed-decider}
  A \emph{typed decider} is a $7$-input Turing machine $\decider$ that on all
  inputs of the form $(n,\lvar,x,\rvar,y,a,b)$ where $n$ is an integer and
  $\lvar,x,\rvar,y,a,b \in \{0,1\}^*$, $\decider$ halts and returns a single
  bit.
  When $\decider$ returns $0$ we say that it \emph{rejects}, otherwise we say
  that it \emph{accepts}.
  We use $\TIME_\decider(n)$ to denote the time complexity of $\decider$ on
  inputs of the form $(n,\ldots)$.
\end{definition}

\begin{definition}
  \label{def:typed-normal-ver}
  Let $\type$ be a finite set and let $G = (\type, E)$ be a graph.
  A $(\type, G)$-\emph{typed normal form verifier} is a pair $\verifier =
  (\sampler, \decider)$ where $\sampler$ is a $(\type,G)$-typed sampler with field
  size $q(n) = 2$ and $\decider$ is a typed decider.
\end{definition}

\begin{definition}
  \label{def:typed-normal-game}
  Let $\verifier = (\sampler, \decider)$ be a $(\type, G)$-typed normal form
  verifier.
  For $n \in \N$, we define the following nonlocal game $\verifier_{n}$ to be
  the \emph{$n$-th game corresponding to the verifier $\verifier$}.
  The question sets $\cal{X}$ and $\cal{Y}$ are $\type \times
  \{0,1\}^{\TIME_\sampler(n)}$.
  The answer sets $\cal{A}$ and $\cal{B}$ are $\{0,1\}^{\TIME_\decider(n)}$.
  The question distribution is the distribution $\mu^G_{\sampler,\, n}$
  specified in Definition~\ref{def:typed-sampler-sample}.
  The decision predicate is the function computed by
  $\decider(n,\lvar,x,\rvar,y,a,b)$, when the inputs $((\lvar, x), (\rvar, y), a, b)$ are
  restricted to $\cal{X}\times \cal{Y}\times\cal{A}\times\cal{B}$.
  The value of the game is denoted by $\val^*(\verifier_{n})$.  
\end{definition}

For $w\in\{\alice,\bob\}$ and a question $(\lvar, x)$ to player $w$ we refer to
$\lvar$ as the \emph{question type} and $x$ as the \emph{question content}.

\subsection{Graph distributions}
\label{sec:graph-dist}

We describe a construction of conditionally linear distributions which sample
from the graph distribution (see Definition~\ref{def:graph-distribution}) of a
graph $G = (U, E)$.
We begin with a technical definition, followed by the definition of the
conditionally linear distribution.

\begin{definition}[Neighbor indicator]\label{def:type-neighbor}
  Given a graph $G = (U, E)$, the \emph{neighbor indicator} of a vertex $u \in
  U$ is the vector $\mathrm{neigh}_G(u) \in \F_2^U$ in which, for all $v \in U$,
  \begin{equation*}
    \mathrm{neigh}_G(u)_v =
    \begin{cases}
      1 & \text{ if } \{u, v\} \in E,\\
      0 & \text{ otherwise}.
    \end{cases}
  \end{equation*}
  In addition, the \emph{$\F_2$-encoding} of a vertex $u \in U$ is the vector
  $\mathrm{enc}_G(u) \in \F_2^U \times \F_2^U$ given by $\mathrm{enc}_G(u) =
  (e_u, \mathrm{neigh}_G(u))$, where $e_u$ is the standard basis vector with
  a~$1$ in the $u$-th position and $0$'s everywhere else.
\end{definition}

\begin{definition}[Graph sampler]\label{def:graph-sampler}
  Let $G = (U, E)$ be a graph with $n$ vertices.
  Then the \emph{conditionally linear functions corresponding to~$G$} are the
  pair of functions $L_G^\alice, L_G^\bob$ on linear space $V_G$ specified in
  \cref{fig:graph-distribution} where $V_G = V_{\vertex{\alice}} \oplus
  V_{\edge{\alice}} \oplus V_{\vertex{\bob}} \oplus V_{\edge{\bob}}$.
\end{definition}

\begin{figure}[!htb]
  \begin{gamespec}
    \centering
    \newcolumntype{C}{>{\centering\arraybackslash $}X<{$}}
    \begin{tabularx}{\textwidth}{*{4}C}
      \multicolumn{4}{c}{\textbf{Subspaces}}\\
      \midrule
      V_{\vertex{\alice}} & V_{\edge{\alice}} & V_{\vertex{\bob}} &
       V_{\edge{\bob}} \\
      \midrule
      \F_2^U & \F_2^U & \F_2^U & \F_2^U\\
      \midrule\\
    \end{tabularx}

    \begin{tabularx}{\textwidth}{p{3.8cm}X}
      \multicolumn{2}{c}{\textbf{Conditionally linear function $L_G^\alice$}}\\
      \midrule
      1st factor subspace & $V_{\vertex{\alice}} \oplus V_{\edge{\alice}}$\\
      1st linear function & Identity function\\
      \midrule
      2nd factor subspace & $V_{\vertex{\bob}} \oplus V_{\edge{\bob}}$\\
      2nd linear functions & For all $x \in V_{\vertex{\alice}} \oplus
      V_{\edge{\alice}}$, suppose there exists a $u \in U$ such that $x =
      \mathrm{enc}_G(u)$.
      Then for all $y \in V_{\vertex{\bob}} \oplus V_{\edge{\bob}}$,
      $L^\alice_{G,\, 2,\, x}$ zeroes out all entries of~$y$ except for
      $(y^{V_{\edge{\bob}}})_u$.
      Otherwise, $L^\alice_{G,\, 2,\, x} = 0$.\\
      \midrule\\
    \end{tabularx}

    \begin{tabularx}{\textwidth}{p{3.8cm}X}
      \multicolumn{2}{c}{\textbf{Conditionally linear function $L^\bob_G$}}\\
      \midrule
      1st factor subspace & $V_{\vertex{\bob}} \oplus V_{\edge{\bob}}$\\
      1st linear function & Identity function\\
      \midrule
      2nd factor subspace & $V_{\vertex{\alice}} \oplus V_{\edge{\alice}}$\\
      2nd linear functions & Similarly defined as those for $L^\alice_G$ by swapping
      $V_{\vertex{\alice}}$ and $V_{\edge{\alice}}$ with $V_{\vertex{\bob}}$ and
      $V_{\edge{\bob}}$ respectively.\\
      \midrule
    \end{tabularx}
  \end{gamespec}
  \caption{Specification of the conditionally linear functions corresponding
    to~$G$.}
  \label{fig:graph-distribution}
\end{figure}

These conditionally linear functions do not simulate the graph distribution
in the sense of sampling directly from it.
The following proposition, however, does show a sense in which these functions
simulate the graph distribution, namely via rejection sampling.

\begin{proposition}[Simulating the graph distribution]
  \label{prop:simulating-graph}
  Let $G = (U, E)$ be a graph with $n$ vertices and $m$ edges, $k$ of which
  are self-loops.
  Let $L^\alice_G$, $L^\bob_G$ be the conditionally linear functions
  corresponding to $G$ (see Figure~\ref{fig:graph-distribution} for the
  definition and associated notation).
  Let $(x, y) \sim \mu_{L^\alice_G,\, L^\bob_G}$.
  Consider the event $\mathcal{E}_G$ that there exists $u, v \in U$ such that
  the following two statements are true.
  \begin{enumerate}[label=(\roman*)]
  \item $x^{V_{\vertex{\alice}} \oplus V_{\edge{\alice}}} = \mathrm{enc}_G(u)$
   	and $y^{V_{\vertex{\bob}} \oplus V_{\edge{\bob}}} = \mathrm{enc}_G(v)$,
  \item  $(x^{V_{\edge{\bob}}})_u = (y^{V_{\edge{\alice}}})_v = 1$.
  \end{enumerate}
  Then
  \begin{enumerate}
  \item \label{item:graph-prob} $\Pr_{x, y} (\mathcal{E}_G) = (2m-k)/16^n$.
  \item \label{item:edge-dist} Conditioned on $\mathcal{E}_G$, $(u, v)$ are
    distributed as the graph distribution of~$G$ (see
    Definition~\ref{def:graph-distribution}).
  \end{enumerate}
  Note that $\mathcal{E}_G$ occurs if and only if both $x^{V_\edge{\bob}}$ and
  $y^{V_\edge{\alice}}$ are nonzero.
  In particular, if $x$ and $y$ are sampled from $\mu_{L^\alice_G,\, L^\bob_G}$
  and given to the respective players, then at least one of them knows when the
  event $\mathcal{E}_G$ does \emph{not} occur.
\end{proposition}

\begin{proof}
  Let $z$ be drawn uniformly at random from $V_{\vertex{\alice}}\oplus
  V_{\edge{\alice}}\oplus V_{\vertex{\bob}} \oplus V_{\edge{\bob}}$, and let $x
  = L^\alice_G(z)$ and $y = L^\bob_G(z)$.
  Then with probability $n^2/16^n$, there exist $u, v \in U$ such that
  \begin{equation*}
    \text{$x^{V_{\vertex{\alice}} \oplus V_{\edge{\alice}}} = \mathrm{enc}_G(u)$
      and $y^{V_{\vertex{\bob}} \oplus V_{\edge{\bob}}} = \mathrm{enc}_G(v)$}\;.
  \end{equation*}
  Conditioned on this occurring, $u$ and~$v$ are distributed as independent,
  uniformly random vertices in~$U$.
  If we further condition on $\{u, v\} \in E$, which occurs with probability
  $(2m-k)/n^2$, then by definition, $(u, v)$ is distributed as the graph
  distribution of~$G$.
  But this event is exactly the event that $\mathcal{E}_G$ holds on~$(x, y)$,
  establishing the proposition.
\end{proof}

\subsection{Detyping typed verifiers}
\label{sec:detype}  

We give a canonical method for taking a typed normal form verifier and producing
an untyped normal form verifier which simulates it.
Throughout this section, $\type$ denotes a finite set, $G = (\type, E)$ denotes
a graph, and $L_G^\alice, L_G^\bob$ denote the conditionally linear functions
corresponding to~$G$ acting on the vector space $V_G = V_{\vertex{\alice}}
\oplus V_{\edge{\alice}} \oplus V_{\vertex{\bob}} \oplus V_{\edge{\bob}}$ of
dimension $4\cdot |\type|$ over $\F_2$, as in
Definition~\ref{def:graph-sampler}.

\begin{definition}[Detyped CL functions]
  \label{def:detyped-CL}
  Let $L^\alice = \{L^{\alice}_{\tvar}\}, L^\bob = \{L^{\bob}_{\tvar}\}$ be
  $\type$-typed families of $\ell$-level conditionally linear functions on $V$.
  We define the \emph{detyped CL functions corresponding to $(L^\alice, L^\bob)$
    on~$G$} to be the pair of $(\ell+2)$-level CL functions $(R^\alice, R^\bob)
  = \detype_G(L^\alice, L^\bob)$ on linear space $V_{\gamestyle{detype}} = V_G
  \oplus V$ as follows.
  For $w \in \{\alice, \bob\}$, and $z \in L^w_G (V_G)$, define the family of
  $\ell$-level CL functions $\{L^w_z\}$ on $V$ as
  \begin{equation*}
    L^w_z =
    \begin{cases}
      0 & \text{ if } z^{V_{\edge{}\overline{w}}} = 0,\\
      L^w_\tvar & \text { otherwise, for } z^{V_{\vertex{}w}} = e_\tvar.
    \end{cases}
  \end{equation*}
  We note that when $z^{V_{\edge{}\overline{w}}}$ is nonzero, it is always the
  case that $z^{V_{\vertex{}w}} = e_\tvar$ for a unique type $\tvar$, by
  Definition~\ref{def:graph-sampler}.
  For $w\in \{\alice, \bob\}$, $R^w$ is the concatenation of $L^w_G$ and
  $\{L^w_z\}_z$ (cf.\ Lemma~\ref{lem:cl-concat}).
\end{definition}

\begin{definition}[Detyped samplers]
  Let $\sampler$ be a $(\type,G)$-typed sampler.
  For each $n \in \N$, let $\{L^{\alice,\, n}_\tvar\}, \{L^{\bob,\, n}_\tvar\}$
  be the CL functions of~$\sampler$ on index~$n$, and set
  $(R^{\alice,\, n}, R^{\bob,\, n}) = \detype_G(L^{\alice,\, n}, L^{\bob,\,
    n})$.
  Then the \emph{detyped sampler} $\detype(\sampler)$ is the (standard)
  sampler whose CL functions on index~$n$ are $R^{\alice,\, n}, R^{\bob,\, n}$.
  Its dimension function is $s_{\gamestyle{detype}}(n) = 4|\type|+s(n)$.
\end{definition}

We note that $\detype(\sampler)$ is indeed a sampler in line with Definition~\ref{def:sampler},
i.e.\ a sampler without types.

\begin{definition}[Detyped deciders]
  Let $\decider$ be a typed decider.
  We define the \emph{detyped decider} $\detype_G(\decider)$ to be the
  (standard) decider that behaves as follows: on input $(n,x,y,a,b)$, it
  attempts to parse $x = (x', x'')$, $y =(y', y'') \in V_G \times \{0,1\}^*$
  (using a canonical scheme for representing pairs of strings).
  If it cannot, it accepts.
  Otherwise, suppose that there exists $\{\lvar, \rvar\} \in E_G$ such that,
  using notation from Definition~\ref{def:type-neighbor},
  \begin{equation*}
    x' = (e_\lvar, \mathrm{neigh}_G(\lvar), 0, e_\lvar)\;,\quad
    y' = (0, e_\rvar, e_\rvar, \mathrm{neigh}_G(\rvar))
    \, \in V_{\vertex{\alice}} \oplus V_{\edge{\alice}}
    \oplus V_{\vertex{\bob}} \oplus V_{\edge{\bob}}\;.
  \end{equation*}
  Then it returns the output of $\decider$ on input $(n, \lvar, x'', \rvar, y'',
  a, b)$.
  Otherwise, it accepts.
\end{definition}

We note that we include the subscript ``$G$'' in $\detype_G(\decider)$
because we do not specify a graph when defining a typed decider (cf.\ Definition~\ref{def:typed-decider}).
This is different than the situation for samplers $\sampler$,
where we do not include a subscript in $\detype(\sampler)$
because the graph $G$ is already specified when defining $\sampler$.

\begin{definition}[Detyped verifiers]
\label{def:detyped-verifier}
  Let $\verifier = (\sampler, \decider)$ be a $(\type, G)$-typed normal form
  verifier.
  We define the \emph{detyped verifier}, denoted by $\detype(\verifier)$, to be
  the (standard) normal form verifier
  $(\detype(\sampler), \detype_G(\decider))$.
\end{definition}

\begin{lemma}[Typed verifiers to detyped verifiers]
  \label{lem:detyping-verifiers}
  Let $\verifier = (\sampler,\decider)$ be a $(\type,G)$-typed normal form
  verifier.
  The detyped verifier $\detype(\verifier) =
  (\detype(\sampler),\detype_G(\decider))$ satisfies the following properties:
  for all $n \in \N$,
  \begin{enumerate}
	\item (\textbf{Completeness}) If $\verifier_n$ has a value-$1$ PCC strategy,
    then $\detype(\verifier)_n$ has a value-$1$ PCC strategy.
	\item (\textbf{Soundness}) If $\val^*(\detype(\verifier)_n) \geq 1 - \eps$,
    then $\val^*(\verifier_n) \geq 1 - 16^{|\type|} \cdot \eps$.
    Furthermore,
    \[
      \Ent(\detype(\verifier)_n,1 - \eps) \geq
      \Ent(\verifier_n,1 - 16^{|\type|} \cdot\eps).
    \]
	\item (\textbf{Sampler parameters}) If $\sampler$ is an $\ell$-level sampler,
    then $\detype(\sampler)$ is an $(\ell+2)$-level sampler.
    The time complexity of $\detype(\sampler)$ satisfies the
    following:
    \begin{gather*}
      \TIME_{\detype(\sampler)}(n) = \poly(|\type|,\TIME_\sampler(n)).
    \end{gather*}
	\item (\textbf{Decider complexity}) The decider $\detype_G(\decider)$ has time
    complexity $\poly(|\type|,\TIME_\decider(n))$.
	\item (\textbf{Efficient computability}) The descriptions of
    $\detype(\sampler)$ and $\detype_G(\decider)$ are polynomial time
    computable from the description of~$G$ and the descriptions of $\sampler$
    and $\decider$, respectively.
  \end{enumerate}
\end{lemma}

\begin{proof}
  Throughout this proof, we fix an index~$n$.
  Let $s = s(n)$ be the dimension of $\sampler$.
  The ambient space of~$\sampler$ is $V = \F_2^s$ and the ambient space
  of~$\detype(\sampler)$ is $V_{\gamestyle{detype}} = V_G \oplus V$.
  Let $\lvar, \rvar \in \type$.
  For this proof, we introduce the notation
  \begin{equation*}
    \mathrm{view}^\alice(\lvar) =
    (e_\lvar, \mathrm{neigh}_G(\lvar), 0, e_\lvar),\;
    \mathrm{view}^\bob(\rvar) =
    (0, e_\rvar, e_\rvar, \mathrm{neigh}_G(\rvar))
    \in V_{\vertex{\alice}} \oplus V_{\edge{\alice}}
    \oplus V_{\vertex{\bob}} \oplus V_{\edge{\bob}}\;.
\end{equation*}
Supposing that players~\alice and~\bob receive~$x$ and~$y$ in
$V_{\gamestyle{detype}}$, and supposing that $(x, y)$ satisfies event
$\mathcal{E}_G$ from Proposition~\ref{prop:simulating-graph}, then $x^{V_G} =
\mathrm{view}^\alice(\lvar)$ and $y^{V_G} = \mathrm{view}^\bob(\rvar)$ for a unique
$\{\lvar, \rvar\} \in E$.

\paragraph{Completeness.}
Let $\strategy = (\ket{\psi},A, B)$ be a value-$1$ PCC strategy for
$\verifier_n$.
We construct a PCC strategy $\strategy^{\gamestyle{DETYPE}}$ for
$\detype(\verifier)_n$ with value~$1$.
This strategy also uses the state $\ket{\psi}$. 
When a player receives a question, they perform measurements described as
follows.
\begin{description}
\item Player \alice: given $x \in V_{\gamestyle{detype}}$ the player checks if
  for some $\lvar \in \type$, $x^{V_G} = \mathrm{view}^\alice(\lvar)$.
	If so, they perform the
  measurement
  \begin{equation*}
    \Big\{ A^{(\lvar,\, x^V)}_{\, a} \Big\}
  \end{equation*}
  to obtain an outcome~$a$, which they use as their answer.
	If not, they reply with the empty string.
	(This entails performing the measurement whose POVM element corresponding to
  the empty string is the identity matrix.)
\item Player \bob: given $y \in V_{\gamestyle{detype}}$, the player checks if
  for some $\rvar \in \type$, $y^{V_G} = \mathrm{view}^\bob(\rvar)$.
	If so, they perform the measurement
  \begin{equation*}
    \Big\{ B^{(\rvar,\, y^V)}_{\, b} \Big\}
  \end{equation*}
  to obtain an outcome~$b$, which they use as their answer.
	If not, they reply with the empty string.
\end{description}
This strategy is projective and consistent because the only measurements it uses
are those in $\strategy$ and ``trivial" measurements containing the identity
matrix.
Suppose the players receive questions~$x$ and~$y$ such that both $x^{V_G} =
\mathrm{view}^\alice(\lvar)$ and $y^{V_G} = \mathrm{view}^\bob(\rvar)$.
In this case, the questions $(\lvar, x^{V})$ and $(\rvar, y^{V})$ are in the
support of the question distribution of $\verifier$.
As a result, the players succeed with probability~$1$ on these questions, and
their measurements always commute.
For the remaining pairs of questions, the decider $\detype_G(\decider)$ always
accepts, and the measurements always commute by virtue of the fact that at least
one is trivial, i.e.\ containing the identity matrix as a POVM element.

\paragraph{Soundness.}
Let $\strategy = (\ket{\psi}, A, B)$ be a strategy for $\detype(\verifier)_n$
with value $1-\eps$.
Suppose $G$ has~$m$ edges, $k$ of which are self-loops.
For any $(x, y)$ drawn from $\mu_{\detype(\sampler),\, n}$, the decider
$\detype_G(\decider)$ automatically accepts unless $(x^{V_G},y^{V_G})$ satisfies
event $\mathcal{E}_G$ from Proposition~\ref{prop:simulating-graph}, which occurs
with probability $(2m-k)/16^{|\type|}$.
When this happens, $x^{V_G}$ and $y^{V_G}$ are distributed as
$\mathrm{view}^\alice(\lvar)$ and $\mathrm{view}^\bob(\rvar)$, where $(\lvar,
\rvar)$ are distributed as the graph distribution on~$G$.
As a result, conditioned on $\mathcal{E}_G$, the probability that $\strategy$
succeeds on $\detype(\verifier)_n$ is equal to the probability that the strategy
$\strategy' = (\ket{\psi}, A', B')$ succeeds on $\verifier_n$, where
\begin{equation*}
  (A')^{\lvar,\, x^V}_a = A^{(\mathrm{view}^\alice(\lvar), x^V)}_a\;, \quad
  (B')^{\rvar,\, y^V}_b = B^{(\mathrm{view}^\bob(\rvar), y^V)}_b\;.
\end{equation*}
This means that
\begin{align*}
  \val^*(\detype(\verifier)_n, \strategy)
  & = \left(1 - \frac{2m-k}{16^{|\type|}}\right) +
    \frac{2m-k}{16^{|\type|}} \cdot \val^*(\verifier_n, \strategy')\\
  & \leq \left(1 - \frac{1}{16^{|\type|}}\right) +
    \frac{1}{16^{|\type|}} \cdot \val^*(\verifier_n, \strategy')\;.
\end{align*}
Thus, $\strategy'$ has value at least $1-16^{|\type|} \cdot \epsilon$.
This proves the first statement in the soundness.
As for the second, $\strategy$ and $\strategy'$ use the same state $\ket{\psi}$,
and therefore both strategies have the same Schmidt rank, which by definition is
at least $\Ent(\verifier_n,1 - 16^{|\type|} \cdot\eps)$.

\paragraph{Complexity.}
Definition~\ref{def:detyped-CL} implies that $\detype(\sampler)$ is an
$(\ell+2)$-level sampler by \cref{lem:cl-concat} and that
$V_{\gamestyle{detype}}$ has dimension $4|\type| + s$.
The claimed time bounds of $\detype(\sampler)$ and $\detype_G(\decider)$
follow from the fact that these perform simple, $\poly(|\type|)$-time
computations followed by running $\sampler$ and $\decider$ as subroutines.
\end{proof}


 
\section{Classical and Quantum Low-degree Tests}
